\scanpage{20}
\mspage{10}
sur $\tau$ par l'op. donnée, et sur $\operatorname{Hom}_{\mathrm{gr}}(T,\pi)$ en considérant $T$, $\pi$ comme des $\mathcal E$-groupes grâce à $\chi$ et à l'opération de $\mathcal E$ sur $\pi$ par autom intérieurs\ldots Pour
\deleted{\deleted{ceci on introduit}
\[
J=\tau/\pi
\]
sur lequel $\mathcal E/\pi$ \inserted{qui est $\mathcal G$} opère par passage au quotient, $\tau\longrightarrow J$ devenant donc un $\mathcal E$-morphisme. Considérons alors le groupoïde $\langle T,J\rangle$ défini par
\[
\operatorname{Ob}\langle T,J\rangle=J
\]
\[
\operatorname{Isom}(j,j')=\left\{\begin{array}{ll}
\varnothing&\text{si }j\ne j'\\
T&\text{si }j=j'
\end{array}\right.
\]
composition des flèches étant celle de $T$. On a une opération de $\mathcal G$ sur $\langle T,J\rangle$ via ses opérations sur $T$ \inserted{(par $\chi$)} et sur $J$. On va définir un $\mathcal G$-morphisme de groupoïdes
\diagram{\begin{tikzcd}[ampersand replacement=\&]
\langle T,J\rangle \arrow[r] \& \Pi^\wedge
\end{tikzcd}}
où $\Pi^\wedge$ est le groupoïde des rev. \inserted{\unsure{0-conn.}} universels de $\Pi$, identifiable (via $\widetilde\Pi$ qui est un \unsure{$\in\Pi^\wedge$}) au groupoïde des}
\mspage{19}
